\noindent\textbf{Problem Statement}

The Segment Anything Model (SAM) has established a new baseline for static image segmentation; however, it is structurally ill-equipped for Referring Audio-Visual Segmentation (Ref-AVS). Current foundation models like SAM suffer from two critical limitations in this context:

\begin{enumerate}
    \item \textbf{Lack of Temporal Awareness:} SAM processes inputs as isolated static frames, failing to capture the temporal consistency and dynamic context necessary for video segmentation.
    \item \textbf{Reliance on Explicit Interaction:} SAM depends on manual user prompts (points, boxes, or masks) to identify targets. It lacks the native ability to interpret implicit ``multimodal prompts''---such as identifying an object described by a specific sound or textual description---without human intervention.
\end{enumerate}

\vspace{1em}
\noindent\textbf{Core Hypothesis}

We hypothesize that we can adapt the frozen, pre-trained knowledge of SAM for dynamic audio-visual scenes without heavy retraining by introducing two novel architectural components. First, a parallel \textbf{temporal modeling branch} can inject time-series context into the static encoder. Second, we can replace the manual ``user-in-the-loop'' prompting mechanism with \textbf{data-driven multimodal prompts}. By translating aligned audio and text features into the sparse and dense prompt formats that SAM natively understands, we believe the model can learn to ``prompt itself'' to segment temporally evolving objects described by complex semantic cues.

\vspace{1em}
\noindent\textbf{Proposed Methodology (High-Level Technical Approach)}

We propose \textbf{TSAM}, a modified architecture that wraps the standard SAM backbone. Our approach focuses on minimal trainable additions while keeping the heavy image encoder frozen.

\vspace{0.5em}
\noindent\textbf{1. Temporal Modeling Branch (Context Injection)}

Instead of retraining the heavyweight image encoder, we will introduce a lightweight auxiliary branch running in parallel.
\begin{itemize}
    \item \textbf{Sequential Processing:} As the frozen encoder processes a frame, this parallel branch will accept features from current and previous steps.
    \item \textbf{Cached Memory \& Adapters:} We will utilize a memory mechanism to store a running summary of previous frame representations. An adapter module will fuse these historical visual states with current audio-text embeddings. This allows the model to maintain object permanence and consistency across time.
\end{itemize}

\vspace{0.5em}
\noindent\textbf{2. Automated Multimodal Prompting}

We aim to synthesize the ``prompts'' SAM expects (points and masks) using audio-visual-text correlations rather than manual clicks.
\begin{itemize}
    \item \textbf{Sparse Prompting Module (Simulating Points):} We will design a query selection mechanism. By analyzing the correlation between the reference text and the audio stream, the system will select the most relevant ``audio cues.'' These cues will be projected and treated as sparse queries (analogous to point prompts) to guide the mask decoder toward the sound source.
    \item \textbf{Dense Prompting Module (Simulating Masks):} We will implement a deep fusion module that uses cross-attention to mix visual features with aligned audio and text embeddings. This will generate a dense, spatial feature map that acts as a coarse ``mask prompt,'' providing global context to the decoder.
\end{itemize}

\vspace{0.5em}
\noindent\textbf{3. Decoding}

The final segmentation will be generated by the standard SAM mask decoder, which will be queried by our synthetic sparse and dense prompts, refined by the temporally aware features from our auxiliary branch.

\vspace{1em}
\noindent\textbf{Expected Contribution}

\begin{itemize}
    \item \textbf{Architectural Novelty:} A framework for extending static, unimodal foundation models (like SAM) into the temporal and multimodal domains without compromising their zero-shot capabilities.
    \item \textbf{Methodological Advancement:} A demonstration of ``latent prompting,'' where semantic cues (audio/text) are mathematically translated into the geometric prompts (points/masks) required by segmentation models.
    \item \textbf{Theoretical Impact:} Establishing that lightweight temporal adapters and cross-modal attention mechanisms are sufficient to unlock video understanding in static image backbones.
\end{itemize}